% Options for packages loaded elsewhere
% Options for packages loaded elsewhere
\PassOptionsToPackage{unicode}{hyperref}
\PassOptionsToPackage{hyphens}{url}
\PassOptionsToPackage{dvipsnames,svgnames,x11names}{xcolor}
%
\documentclass[
]{report}
\usepackage{xcolor}
\usepackage[top=30mm,left=20mm,right=20mm,bottom=30mm]{geometry}
\usepackage{amsmath,amssymb}
\setcounter{secnumdepth}{5}
\usepackage{iftex}
\ifPDFTeX
  \usepackage[T1]{fontenc}
  \usepackage[utf8]{inputenc}
  \usepackage{textcomp} % provide euro and other symbols
\else % if luatex or xetex
  \usepackage{unicode-math} % this also loads fontspec
  \defaultfontfeatures{Scale=MatchLowercase}
  \defaultfontfeatures[\rmfamily]{Ligatures=TeX,Scale=1}
\fi
\usepackage{lmodern}
\ifPDFTeX\else
  % xetex/luatex font selection
\fi
% Use upquote if available, for straight quotes in verbatim environments
\IfFileExists{upquote.sty}{\usepackage{upquote}}{}
\IfFileExists{microtype.sty}{% use microtype if available
  \usepackage[]{microtype}
  \UseMicrotypeSet[protrusion]{basicmath} % disable protrusion for tt fonts
}{}
\makeatletter
\@ifundefined{KOMAClassName}{% if non-KOMA class
  \IfFileExists{parskip.sty}{%
    \usepackage{parskip}
  }{% else
    \setlength{\parindent}{0pt}
    \setlength{\parskip}{6pt plus 2pt minus 1pt}}
}{% if KOMA class
  \KOMAoptions{parskip=half}}
\makeatother
% Make \paragraph and \subparagraph free-standing
\makeatletter
\ifx\paragraph\undefined\else
  \let\oldparagraph\paragraph
  \renewcommand{\paragraph}{
    \@ifstar
      \xxxParagraphStar
      \xxxParagraphNoStar
  }
  \newcommand{\xxxParagraphStar}[1]{\oldparagraph*{#1}\mbox{}}
  \newcommand{\xxxParagraphNoStar}[1]{\oldparagraph{#1}\mbox{}}
\fi
\ifx\subparagraph\undefined\else
  \let\oldsubparagraph\subparagraph
  \renewcommand{\subparagraph}{
    \@ifstar
      \xxxSubParagraphStar
      \xxxSubParagraphNoStar
  }
  \newcommand{\xxxSubParagraphStar}[1]{\oldsubparagraph*{#1}\mbox{}}
  \newcommand{\xxxSubParagraphNoStar}[1]{\oldsubparagraph{#1}\mbox{}}
\fi
\makeatother


\usepackage{longtable,booktabs,array}
\usepackage{calc} % for calculating minipage widths
% Correct order of tables after \paragraph or \subparagraph
\usepackage{etoolbox}
\makeatletter
\patchcmd\longtable{\par}{\if@noskipsec\mbox{}\fi\par}{}{}
\makeatother
% Allow footnotes in longtable head/foot
\IfFileExists{footnotehyper.sty}{\usepackage{footnotehyper}}{\usepackage{footnote}}
\makesavenoteenv{longtable}
\usepackage{graphicx}
\makeatletter
\newsavebox\pandoc@box
\newcommand*\pandocbounded[1]{% scales image to fit in text height/width
  \sbox\pandoc@box{#1}%
  \Gscale@div\@tempa{\textheight}{\dimexpr\ht\pandoc@box+\dp\pandoc@box\relax}%
  \Gscale@div\@tempb{\linewidth}{\wd\pandoc@box}%
  \ifdim\@tempb\p@<\@tempa\p@\let\@tempa\@tempb\fi% select the smaller of both
  \ifdim\@tempa\p@<\p@\scalebox{\@tempa}{\usebox\pandoc@box}%
  \else\usebox{\pandoc@box}%
  \fi%
}
% Set default figure placement to htbp
\def\fps@figure{htbp}
\makeatother





\setlength{\emergencystretch}{3em} % prevent overfull lines

\providecommand{\tightlist}{%
  \setlength{\itemsep}{0pt}\setlength{\parskip}{0pt}}



 


\makeatletter
\@ifpackageloaded{caption}{}{\usepackage{caption}}
\AtBeginDocument{%
\ifdefined\contentsname
  \renewcommand*\contentsname{Table of contents}
\else
  \newcommand\contentsname{Table of contents}
\fi
\ifdefined\listfigurename
  \renewcommand*\listfigurename{List of Figures}
\else
  \newcommand\listfigurename{List of Figures}
\fi
\ifdefined\listtablename
  \renewcommand*\listtablename{List of Tables}
\else
  \newcommand\listtablename{List of Tables}
\fi
\ifdefined\figurename
  \renewcommand*\figurename{Figure}
\else
  \newcommand\figurename{Figure}
\fi
\ifdefined\tablename
  \renewcommand*\tablename{Table}
\else
  \newcommand\tablename{Table}
\fi
}
\@ifpackageloaded{float}{}{\usepackage{float}}
\floatstyle{ruled}
\@ifundefined{c@chapter}{\newfloat{codelisting}{h}{lop}}{\newfloat{codelisting}{h}{lop}[chapter]}
\floatname{codelisting}{Listing}
\newcommand*\listoflistings{\listof{codelisting}{List of Listings}}
\makeatother
\makeatletter
\makeatother
\makeatletter
\@ifpackageloaded{caption}{}{\usepackage{caption}}
\@ifpackageloaded{subcaption}{}{\usepackage{subcaption}}
\makeatother
\usepackage{bookmark}
\IfFileExists{xurl.sty}{\usepackage{xurl}}{} % add URL line breaks if available
\urlstyle{same}
\hypersetup{
  colorlinks=true,
  linkcolor={blue},
  filecolor={Maroon},
  citecolor={Blue},
  urlcolor={Blue},
  pdfcreator={LaTeX via pandoc}}


\author{}
\date{}
\begin{document}

\renewcommand*\contentsname{Table of contents}
{
\hypersetup{linkcolor=}
\setcounter{tocdepth}{2}
\tableofcontents
}

\chapter{AI Capability Checkpoints}\label{ai-capability-checkpoints}

This document provides optional reflection scaffolds aligned to a
six-domain AI Capability model.\\
These checkpoints are designed to support responsible, transparent, and
professional research practice.

Completion is \textbf{not required}.\\
If AI tools were not used in your research process, this section may be
left blank or removed.

\begin{center}\rule{0.5\linewidth}{0.5pt}\end{center}

\section{1. Awareness \& Orientation}\label{awareness-orientation}

If AI tools were used at this stage of the project:

\begin{itemize}
\tightlist
\item
  What tools were used (e.g., drafting assistants, coding support,
  summarisation tools)?
\item
  What are their known limitations (e.g., hallucination, bias
  amplification, incomplete context)?
\item
  What checks were performed to ensure accuracy and reliability?
\end{itemize}

If no AI tools were used, briefly state that the work was completed
without AI assistance.

\begin{center}\rule{0.5\linewidth}{0.5pt}\end{center}

\section{2. Human--AI Co-Agency}\label{humanai-co-agency}

If applicable:

\begin{itemize}
\tightlist
\item
  Which tasks were supported by AI tools?
\item
  Which decisions remained fully under human judgement?
\item
  Where was human oversight explicitly exercised?
\item
  How were AI-generated suggestions evaluated or rejected?
\end{itemize}

Clarify that ultimate responsibility for the work remains with the
research team.

\begin{center}\rule{0.5\linewidth}{0.5pt}\end{center}

\section{3. Applied Practice \&
Innovation}\label{applied-practice-innovation}

If AI tools were used:

\begin{itemize}
\tightlist
\item
  In what specific ways did AI support drafting, structuring, analysis,
  or critique?
\item
  What elements of the research were developed independently?
\item
  Were any AI outputs modified substantially before inclusion?
\end{itemize}

Avoid overstating AI contributions. Emphasise verification and
professional judgement.

\begin{center}\rule{0.5\linewidth}{0.5pt}\end{center}

\section{4. Ethics, Equity \& Impact}\label{ethics-equity-impact}

If AI tools were involved:

\begin{itemize}
\tightlist
\item
  Were any sensitive or identifiable data shared with AI systems?
\item
  What bias or fairness risks were considered?
\item
  Could AI-assisted language or framing unintentionally misrepresent
  findings?
\item
  How were equity implications reviewed?
\end{itemize}

If AI was not used, consider whether the research topic itself carries
equity or impact considerations.

\begin{center}\rule{0.5\linewidth}{0.5pt}\end{center}

\section{5. Decision-Making \&
Governance}\label{decision-making-governance}

If AI tools were used:

\begin{itemize}
\tightlist
\item
  How was AI use documented (e.g., in decision logs or methods notes)?
\item
  Were institutional, funder, or journal disclosure requirements
  considered?
\item
  How were AI-assisted decisions justified and recorded?
\end{itemize}

Ensure transparency and traceability where relevant.

\begin{center}\rule{0.5\linewidth}{0.5pt}\end{center}

\section{6. Reflection, Learning \&
Renewal}\label{reflection-learning-renewal}

Optional reflective prompts:

\begin{itemize}
\tightlist
\item
  What worked well (or not) when integrating AI into this project?
\item
  What would you do differently in future work?
\item
  Did AI influence the framing, clarity, or structure of the research?
\end{itemize}

If AI was not used, reflect briefly on whether it would have added value
or risk in this context.

\begin{center}\rule{0.5\linewidth}{0.5pt}\end{center}

\section{Important Notes}\label{important-notes}

\begin{itemize}
\tightlist
\item
  AI use is \textbf{not required} for this template.
\item
  These checkpoints are designed to support transparency and governance.
\item
  Researchers may adapt, shorten, or remove this section to suit their
  workflow.
\item
  Responsibility for all research outputs remains with the human
  authors.
\end{itemize}




\end{document}
